\documentclass[12pt,oneside]{book}

% ============================================================
% METADATOS
% ============================================================
\newcommand{\universidad}{Universidad de Buenos Aires (UBA)}
\newcommand{\facultad}{Facultad de Derecho}
\newcommand{\carrera}{Consultorio Jurídico Gratuito}
\newcommand{\materia}{Práctica Profesional}
\newcommand{\titulo}{Patrocinio Jurídico Gratuito: derivaciones, cierre de expedientes, bajas administrativas y caso práctico de alimentos}
\newcommand{\autor}{Isaac Edgar Camacho Ocampo}
\newcommand{\anio}{2026}

% ============================================================
% PAQUETES
% ============================================================
\usepackage[utf8]{inputenc}
\usepackage[T1]{fontenc}
\usepackage[spanish,es-nodecimaldot]{babel}
\usepackage[
    top=2.5cm,
    bottom=2.5cm,
    left=3cm,
    right=2.5cm,
    headheight=15pt
]{geometry}
\usepackage{amsmath}
\usepackage{amssymb}
\usepackage{mathtools}
\usepackage{graphicx}
\usepackage{float}
\usepackage{caption}
\usepackage{subcaption}
\usepackage{booktabs}
\usepackage{array}
\usepackage{longtable}
\usepackage{tabularx}
\usepackage{enumitem}
\usepackage{xcolor}
\usepackage[most]{tcolorbox}
\usepackage{fancyhdr}
\usepackage{ragged2e}
\usepackage[
    colorlinks=true,
    linkcolor=blue!50!black,
    citecolor=green!40!black,
    urlcolor=blue!60!black,
    pdftitle={\titulo},
    pdfauthor={\autor},
    pdfsubject={Apuntes universitarios},
    bookmarks=true
]{hyperref}

% ============================================================
% CONFIGURACIÓN
% ============================================================
\graphicspath{
    {./img/}
    {./graficos/}
    {./figuras/}
}
\setcounter{tocdepth}{2}
\setlength{\parindent}{0pt}
\setlength{\parskip}{0.5em}
\setlist[itemize]{
    topsep=0.3em,
    itemsep=0.2em
}
\setlist[enumerate]{
    topsep=0.3em,
    itemsep=0.2em
}
\clubpenalty=10000
\widowpenalty=10000

% ============================================================
% COMANDOS
% ============================================================
\newcommand{\abs}[1]{\left|#1\right|}
\newcommand{\norm}[1]{\left\lVert#1\right\rVert}
\newcommand{\vect}[1]{\mathbf{#1}}
\newcommand{\deriv}[2]{\frac{d#1}{d#2}}
\newcommand{\pderiv}[2]{\frac{\partial#1}{\partial#2}}

% ============================================================
% ENTORNOS / CAJAS
% ============================================================
\newtcolorbox{definition}[1]{
    enhanced,
    breakable,
    title={Definición: #1},
    fonttitle=\bfseries,
    colback=blue!3,
    colframe=blue!50!black,
    boxrule=0.8pt,
    arc=2mm
}

\newtcolorbox{important}{
    enhanced,
    breakable,
    title={Importante},
    fonttitle=\bfseries,
    colback=yellow!8,
    colframe=orange!70!black,
    boxrule=0.8pt,
    arc=2mm
}

\newtcolorbox{example}[1]{
    enhanced,
    breakable,
    title={Ejemplo: #1},
    fonttitle=\bfseries,
    colback=green!4,
    colframe=green!40!black,
    boxrule=0.8pt,
    arc=2mm
}

\newtcolorbox{formula}[1]{
    enhanced,
    breakable,
    title={Fórmula / Regla: #1},
    fonttitle=\bfseries,
    colback=purple!4,
    colframe=purple!50!black,
    boxrule=0.8pt,
    arc=2mm
}

\newtcolorbox{exercise}[1]{
    enhanced,
    breakable,
    title={Ejercicio: #1},
    fonttitle=\bfseries,
    colback=gray!5,
    colframe=gray!60!black,
    boxrule=0.8pt,
    arc=2mm
}

\newtcolorbox{note}{
    enhanced,
    breakable,
    title={Nota},
    fonttitle=\bfseries,
    colback=white,
    colframe=gray!50,
    boxrule=0.6pt,
    arc=2mm
}

% ============================================================
% CONFIGURACIÓN FINAL
% ============================================================
\pagestyle{fancy}
\fancyhf{}
\fancyhead[L]{\nouppercase{\leftmark}}
\fancyhead[R]{\materia}
\fancyfoot[C]{\thepage}
\renewcommand{\headrulewidth}{0.4pt}
\renewcommand{\footrulewidth}{0pt}

\fancypagestyle{plain}{
    \fancyhf{}
    \fancyfoot[C]{\thepage}
    \renewcommand{\headrulewidth}{0pt}
}

% ============================================================
% DOCUMENTO
% ============================================================
\begin{document}

% ------------------------------------------------------------
% PORTADA
% ------------------------------------------------------------
\begin{titlepage}

\thispagestyle{empty}

\begin{center}

\vspace*{1cm}

{\large \textsc{\universidad}}

\vspace{0.2cm}

{\large \textsc{\facultad}}

\vspace{0.2cm}

{\large \carrera}

\vspace{3cm}

{\Huge \textbf{\materia}}

\vspace{1cm}

{\LARGE Patrocinio Jurídico Gratuito}

\vspace{0.4cm}

{\large Derivaciones, cierre de expedientes, bajas administrativas y caso práctico de alimentos}

\vspace{0.8cm}

\textit{Comprender los conceptos es el primer paso para convertir el estudio en conocimiento.}

\vspace{1.2cm}

\rule{0.70\textwidth}{0.5pt}

\vspace{0.5cm}

{\large Apuntes de estudio}

\vspace{0.5cm}

\rule{0.70\textwidth}{0.5pt}

\vfill

{\large \autor}

\vspace{0.3cm}

{\large \anio}

\end{center}

\vfill

\begin{flushleft}
\textit{Este es un modesto aporte a los estudiantes de ``Práctica Profesional'' no pretende sustituir el material oficial de la materia.}
\end{flushleft}

\end{titlepage}

% ------------------------------------------------------------
% ÍNDICE
% ------------------------------------------------------------
\frontmatter
\tableofcontents
\mainmatter

% ------------------------------------------------------------
% INTRODUCCIÓN
% ------------------------------------------------------------
\chapter*{Presentación del material}
\addcontentsline{toc}{chapter}{Presentación del material}
\markboth{Presentación del material}{Presentación del material}

Este material corresponde a una práctica de comisión organizada en tres bloques:

\begin{enumerate}
    \item \textbf{Derivación de consultantes}: cómo se deriva a un consultante que ya no puede seguir en su horario o comisión, con las causas admitidas, la autorización de la coordinadora, la documentación necesaria y el mail <<listo para reenviar>>.
    \item \textbf{Cierre del expediente y bajas administrativas}: qué se hace en el expediente después de derivar (renuncia y pedido de desvinculación) y cómo se cargan las bajas administrativas en el drive.
    \item \textbf{Caso práctico de alimentos}: apelación de los alimentos provisorios, demanda por atribución del hogar, divorcio y la duda sobre si corresponde demandar primero al padre o a la abuela.
\end{enumerate}

\section*{Utilidad para la vida profesional}

El docente aclara que el trámite exacto es propio de la comisión (<<no les va a servir en la vida\ldots{} pero sí sirve mucho acá, aquí y ahora>>). Lo que acompaña a lo largo de la carrera son los hábitos que hay detrás: transferir un caso sin que el cliente quede en el aire, documentar cada paso y cerrar lo que se abre.

\begin{enumerate}
    \item \textbf{Responsabilidad ante el cliente.} Una derivación bien hecha evita que una persona quede sin patrocinio. Un estudio jurídico o un organismo público se maneja igual cuando un colega se va o cambia de área.
    \item \textbf{Protección propia.} Renunciar y pedir la desvinculación evita seguir recibiendo notificaciones de expedientes donde ya no se es parte, un riesgo real de plazos y confusión.
    \item \textbf{Trazabilidad.} Dejar registro (drive, formulario de baja) permite responder dentro de dos años <<qué pasó>> con un cliente. Eso se llama orden, y los empleadores lo valoran.
    \item \textbf{Estrategia procesal.} El caso de alimentos muestra cómo se ordenan los pasos (apelar, demandar la atribución del hogar, divorcio) y el principio de subsidiariedad (primero el progenitor, después la abuela): litigio de familia real.
\end{enumerate}

\section*{Relación entre los bloques}

Los tres bloques pueden compararse con el pase de guardia de un hospital. Cuando un médico termina su turno, no puede dejar al paciente sin atender: lo pasa a un colega con la historia clínica completa (bloque 1: derivación, con declaración jurada, constancia y destino). Después anota en el libro de guardia que el paciente ya no está a su cargo y deja de recibir los avisos de ese paciente (bloque 2: renuncia, desvinculación y baja registrada). Y mientras tanto, cada paciente sigue su tratamiento con un plan por etapas y un orden de prioridades (bloque 3: el caso de alimentos, con sus pasos y su orden de demandados).

Los tres bloques son el mismo trabajo visto desde tres momentos: \textbf{pasar el caso, cerrarlo y llevarlo adelante}.

% ============================================================
% CAPÍTULO 1
% ============================================================
\chapter[Derivación de consultantes]{Derivación de consultantes entre comisiones}

\section{Cuándo corresponde derivar}

Una derivación se plantea cuando un consultante manifiesta que no puede seguir siendo atendido en su comisión. La derivación procede únicamente por una causa real y entendible: principalmente un \textbf{cambio laboral} o una situación extrema, como la de un consultante que lleva a su hijo a terapias por la mañana y por eso se le complica asistir.

El patrocinio jurídico gratuito no depende de que la persona guste o no. Si el motivo no es uno de los dos mencionados, el consultante no puede ser atendido en el consultorio: queda fuera del patrocinio y tampoco puede pasar a otra comisión. El patrocinio es uno solo; por eso no se puede atender a las dos partes.

% TODO: contenido incompleto o ambiguo en las notas originales (el docente califica de <<extremo>> el caso de las terapias; no se aclara si existen otras causas admitidas).

\begin{important}
La derivación solo procede por una causa real y entendible: principalmente un cambio laboral o una situación como llevar a un hijo a terapias en el horario de la comisión. El patrocinio gratuito es uno solo, por lo que no se atiende a las dos partes.
\end{important}

\section{Cambio de horario y autorización de la coordinadora}

Si el consultante tiene una causa real y entendible para cambiar, se puede cambiar de horario. En la facultad hay horarios de 8 de la mañana a 6 de la tarde. Algunos son específicos: los de las 6 de la tarde corresponden a discapacidad y, según el docente, también hay comisiones que trabajan civil y capacidad.

% TODO: contenido incompleto o ambiguo en las notas originales (el docente no recuerda con precisión qué otras materias se atienden en el horario de las 18 h).

El procedimiento es el siguiente:

\begin{enumerate}
    \item Corroborar el motivo del cambio. Ejemplo: el consultante cambió de trabajo y ya no puede asistir a las 8, pero sí a partir de las 14.
    \item Pedir autorización a \textbf{Silvia}, la coordinadora, informándole que la persona puede asistir a partir de esa hora.
\end{enumerate}

\section{Documentación que acredita el motivo}

Para acreditar que el consultante cambió de trabajo se utiliza:

\begin{itemize}
    \item el \textbf{certificado de trabajo}, si el empleo es en blanco (el docente señala que lamentablemente casi nunca lo es); o
    \item una \textbf{nota firmada} por el consultante o la consultante, en la que diga que cambió de trabajo y que trabaja en otro lugar y en otro horario.
\end{itemize}

\begin{note}
En la práctica del caso comentado, un alumno informa que Silvia permitió presentar una nota en lugar del certificado de trabajo, y el docente lo consideró correcto.
\end{note}

\section{El mail <<listo para reenviar>>}

Una vez que Silvia indica a qué comisión debe ir el consultante (por ejemplo, la comisión del doctor García, a las 4 de la tarde), el docente debe avisar a esa comisión: no puede llegar una persona sin aviso.

Para ello, los alumnos deben enviar un correo electrónico a su comisión \textbf{listo para reenviar}. El docente solo debe reenviarlo; el trabajo corresponde a los alumnos, no a la comisión.

La declaración jurada se encuentra en el primer mail que recibió el alumno, que puede ser de hace tres años; hay que buscarla y, si no está, se avisa al docente.

\begin{table}[H]
\centering
\caption{Contenido del mail de derivación}
\begin{tabularx}{\textwidth}{
    >{\raggedright\arraybackslash}p{0.22\textwidth}
    >{\raggedright\arraybackslash}X
    >{\raggedright\arraybackslash}X
}
\toprule
\textbf{Elemento} & \textbf{Dónde se obtiene} & \textbf{Observación del docente} \\
\midrule
\textbf{1. Declaración jurada} &
En el primer mail que recibió el alumno, aunque haya sido hace tres años. &
<<Hay que buscarla, pero está>>. Si no está, se avisa al docente. \\
\addlinespace
\textbf{2. Constancia del motivo} &
Certificado de trabajo, o nota firmada por el consultante o la consultante. &
El certificado casi nunca es <<en blanco>>; por eso suele ser la nota. \\
\addlinespace
\textbf{3. Comisión de destino} &
Indicada por Silvia, la coordinadora (por ejemplo, la comisión del doctor García, a las 4 de la tarde). &
Debe ir la dirección de la comisión para que el docente sepa adónde reenviar. \\
\bottomrule
\end{tabularx}
\end{table}

El docente remarca que son \textbf{tres cosas, no cuatro}: es muy simple, pero parece que todavía no lo han incorporado.

\section{Dudas frecuentes sobre la derivación}

Ante la consulta de un alumno, que informa que en el segundo mail está incorporada la declaración jurada y pregunta si está listo para reenviar, el docente responde que lo único que faltaría es la \textbf{comisión} de destino, y pide que se lo envíen igualmente. Aclara que los alumnos hablan con Silvia y son ellos quienes entregan la información al docente, aunque este a veces también hable con ella por una derivación.

% TODO: contenido incompleto o ambiguo en las notas originales (el primer fragmento de la duda del alumno es ininteligible en la grabación).

\begin{important}
El docente insiste en que no se debe seguir derivando mal: el mail debe llevar siempre declaración jurada, constancia del motivo y comisión de destino.
\end{important}

\subsection*{Cuadro Cornell: derivación de consultantes}

\begin{table}[H]
\centering
\begin{tabularx}{\textwidth}{
    >{\raggedright\arraybackslash}p{0.28\textwidth}
    >{\raggedright\arraybackslash}X
}
\toprule
\textbf{Preguntas / palabras clave} &
\textbf{Notas / respuestas} \\
\midrule

\textbf{¿Cuándo corresponde derivar a un consultante?} &
Solo por una causa real y entendible: principalmente un cambio laboral o, por ejemplo, llevar a un hijo a terapias en el horario de la comisión. Si el motivo no es de esos, el consultante queda fuera del patrocinio. El patrocinio es uno, por lo que no se atiende a las dos partes. \\

\textbf{¿Quién autoriza el cambio de horario o de comisión?} &
Silvia, la coordinadora. Hay horarios de 8 a 18 h; algunos son específicos (los de las 18 h, por ejemplo, de discapacidad). Primero se corrobora el motivo y luego se pide la autorización. \\

\textbf{¿Qué acredita el motivo del cambio?} &
El certificado de trabajo (si es en blanco) o una nota firmada por el consultante o la consultante diciendo que cambió de trabajo y trabaja en otro horario. Silvia aceptó la nota en lugar del certificado. \\

\textbf{¿Qué lleva el mail de derivación?} &
Tres cosas: la declaración jurada (está en el primer mail recibido), la nota o certificado, y la comisión de destino con su dirección. Debe estar <<listo para reenviar>>: el trabajo es del alumno. \\

\midrule

\multicolumn{2}{p{\dimexpr\textwidth-2\tabcolsep\relax}}{
\textbf{Resumen:}
La derivación solo procede por una causa real (típicamente un cambio laboral) y la autoriza la coordinadora. Se acredita con certificado de trabajo o nota firmada y se concreta con un mail listo para reenviar que lleva declaración jurada, constancia del motivo y comisión de destino.
} \\

\bottomrule
\end{tabularx}
\end{table}

% ============================================================
% CAPÍTULO 2
% ============================================================
\chapter[Cierre del expediente y bajas]{Cierre del expediente y bajas administrativas}

\section{Renuncia al patrocinio y pedido de desvinculación}

Después de realizar la derivación hay que tener en cuenta lo que sigue. Se \textbf{renuncia} en el expediente: ya no se tiene nada más que ver allí. Y, lo más importante, se pide la \textbf{desvinculación}.

El docente explica que existen muchos expedientes en los que ya no se es parte y de todos modos siguen apareciendo. Muchas veces se pide la desvinculación y el juzgado no desvincula; aun así, se pide.

\begin{definition}{Desvinculación}
Es el pedido, una vez presentada la renuncia, para que el usuario del consultorio deje de estar vinculado al expediente. Así no se reciben notificaciones de causas en las que ya no se interviene ni se generan confusiones (de lo contrario, además de la confusión, habría que volver a aclararlo).
\end{definition}

\section{Registro de la derivación en el drive}

Una vez que se renuncia y se pide la desvinculación, la derivación se \textbf{registra en el drive} (o donde se prefiera), dejando asentado que fue una derivación.

\section{Bajas administrativas}

Las bajas administrativas son muy importantes. Se realizan desde el drive, mediante un formulario, y tienen estas características:

\begin{itemize}
    \item Deben hacerse también administrativamente para que no queden pendientes: los alumnos se egresan y se van, y al año siguiente puede presentarse un consultante que dice haber estado allí hace dos años sin recordar lo ocurrido. El registro permite responder, por ejemplo, <<te dieron de baja en el 2020>>, y saber qué pasó.
    \item En el drive, las bajas están en una carpeta llamada <<Bajas>>, donde hay un formulario de Google Drive. Es muy simple: tiene cuatro o cinco renglones para completar.
    \item Se hacen porque se acerca el final del cuatrimestre y todo esto tiene que ver con la prolijidad y el orden que se quiere para el \textbf{traspaso}. Si las bajas no están bien cargadas, se puede seguir trabajando en algo que ni se sabe que existe.
\end{itemize}

% TODO: contenido incompleto o ambiguo en las notas originales (fragmento poco claro en 5:12 sobre anotar o fotografiar algo; y fragmento ininteligible entre 6:39 y 7:36 de la grabación).

\begin{table}[H]
\centering
\caption{Secuencia de cierre}
\begin{tabularx}{\textwidth}{
    >{\raggedright\arraybackslash}p{0.25\textwidth}
    >{\raggedright\arraybackslash}X
    >{\raggedright\arraybackslash}X
}
\toprule
\textbf{Paso} & \textbf{Acción} & \textbf{Finalidad} \\
\midrule
\textbf{1. Renuncia} &
Renunciar al patrocinio en el expediente. &
Dejar de intervenir en la causa. \\
\addlinespace
\textbf{2. Desvinculación} &
Pedir al juzgado la desvinculación del usuario. &
No quedar notificados ni vinculados al expediente. \\
\addlinespace
\textbf{3. Registro de la derivación} &
Anotar en el drive que fue una derivación. &
Dejar constancia de lo ocurrido. \\
\addlinespace
\textbf{4. Baja administrativa} &
Completar el formulario de la carpeta <<Bajas>> del drive (cuatro o cinco renglones). &
Saber qué pasó con cada consultante aun años después; ordenar el traspaso. \\
\bottomrule
\end{tabularx}
\end{table}

\subsection*{Cuadro Cornell: cierre del expediente y bajas}

\begin{table}[H]
\centering
\begin{tabularx}{\textwidth}{
    >{\raggedright\arraybackslash}p{0.28\textwidth}
    >{\raggedright\arraybackslash}X
}
\toprule
\textbf{Preguntas / palabras clave} &
\textbf{Notas / respuestas} \\
\midrule

\textbf{¿Qué se hace en el expediente al derivar?} &
Se renuncia al patrocinio y se pide la desvinculación para no quedar notificados. Se registra en el drive que fue una derivación. \\

\textbf{¿Qué es la desvinculación?} &
Que el usuario del consultorio deje de estar vinculado al expediente: no hay nada más que hacer ni notificaciones. El juzgado muchas veces no desvincula, pero se pide igual. \\

\textbf{¿Qué es la baja administrativa y cómo se hace?} &
Un registro interno para saber qué pasó con cada consultante aunque el alumno se egrese. Se completa el formulario de la carpeta <<Bajas>> del drive (cuatro o cinco renglones). Es clave para el traspaso de fin de cuatrimestre. \\

\midrule

\multicolumn{2}{p{\dimexpr\textwidth-2\tabcolsep\relax}}{
\textbf{Resumen:}
Después de derivar se renuncia, se pide la desvinculación para no quedar notificados, se registra la derivación en el drive y se carga la baja administrativa, lo que ordena el traspaso de fin de cuatrimestre.
} \\

\bottomrule
\end{tabularx}
\end{table}

% ============================================================
% CAPÍTULO 3
% ============================================================
\chapter[Caso práctico: alimentos]{Caso práctico: alimentos, atribución del hogar y divorcio}

\section{Relato del caso: apelación de los alimentos provisorios}

Un alumno explica que el pedido de alimentos fue realizado el \textbf{13 de octubre}. Se está preparando un \textbf{recurso de apelación} para cuestionar lo resuelto sobre los alimentos provisorios. Es necesario \textbf{subir primero el escrito apelando}. El docente valora la explicación como <<muy técnica>> y <<muy bien>>.

% TODO: contenido incompleto o ambiguo en las notas originales (el alumno manifiesta no saber cómo explicar el objeto exacto de la apelación; el fragmento está incompleto).

\section{Estrategia procesal: cuota, atribución del hogar y divorcio}

Según lo expuesto por los alumnos, el planteo se ordena en las siguientes etapas: se busca que se fije la cantidad de alimentos, luego se iniciará una demanda por la atribución del hogar (expresada también como <<restitución del hogar>>) para que él pueda volver, y se planteará la decisión de la persona que quiere divorciarse.

% TODO: contenido incompleto o ambiguo en las notas originales (fragmento ininteligible entre 8:51 y 9:57; no queda claro el sujeto que vuelve al hogar ni a qué parte representa el consultorio).

\begin{table}[H]
\centering
\caption{Etapas que plantea el caso}
\begin{tabularx}{\textwidth}{
    >{\centering\arraybackslash}p{0.12\textwidth}
    >{\raggedright\arraybackslash}X
    >{\raggedright\arraybackslash}X
}
\toprule
\textbf{Orden} & \textbf{Actuación} & \textbf{Objetivo} \\
\midrule
\textbf{1} &
Escrito de recurso de apelación (se sube primero). &
Cuestionar lo resuelto sobre los alimentos provisorios, en relación con el pedido del 13 de octubre. \\
\addlinespace
\textbf{2} &
Que se fije la cuota de alimentos. &
Obtener la cantidad de alimentos correspondiente. \\
\addlinespace
\textbf{3} &
Demanda por atribución (<<restitución>>) del hogar. &
Que él pueda volver. \\
\addlinespace
\textbf{4} &
Planteo de divorcio. &
Canalizar la decisión de la persona que quiere divorciarse. \\
\bottomrule
\end{tabularx}
\end{table}

\section{Consultas del alumnado sobre el expediente}

Los alumnos plantearon dudas sobre el expediente. Las intervenciones se conservan de forma fragmentaria, pues la grabación se pierde en varios tramos. Se mencionan:

\begin{itemize}
    \item que ya se apeló y se pretende pedir algo más;
    \item que la persona tenía en \textbf{Paraguay} otros terrenos;
    \item que algo no se subió al expediente;
    \item que \textbf{pueden ser cuatro sujetos};
    \item la referencia a \textbf{alimentando} e \textbf{informe}.
\end{itemize}

% TODO: contenido incompleto o ambiguo en las notas originales (las intervenciones entre 10:22 y 12:57 están cortadas y no permiten reconstruir la consulta completa ni su respuesta).

\section{A quién demandar primero: la obligación subsidiaria de la abuela}

Un alumno consulta si se puede demandar primero al padre o a la abuela, y si el reclamo a la abuela se plantea <<subsidiariamente>>. En la grabación no se escucha el número del artículo al que se alude.

% TODO: contenido incompleto o ambiguo en las notas originales (número de artículo inaudible en 13:12).

El docente explica que la obligación de la abuela es \textbf{subsidiaria} y que hay que empezar por uno de los obligados: no se pasa a la abuela hasta que el progenitor diga que no. El orden es el siguiente:

\begin{enumerate}
    \item \textbf{Mediar contra el progenitor.}
    \item Si el progenitor no se presenta, se cierra el acta y queda una \textbf{constancia}.
    \item Recién entonces se hace la demanda.
    \item Se va contra el progenitor, a demanda y a juicio, y \textbf{después del juicio} se pasa a la abuela.
\end{enumerate}

% TODO: contenido incompleto o ambiguo en las notas originales (el docente dice que quedaría una <<tareita>> por hacer, pero el fragmento es incompleto).

\begin{important}
La obligación de la abuela es subsidiaria: se empieza por el progenitor (mediación; si no se presenta se cierra el acta y queda una constancia; demanda y juicio) y recién después se pasa a la abuela.
\end{important}

\begin{table}[H]
\centering
\caption{Orden de los demandados según lo explicado}
\begin{tabularx}{\textwidth}{
    >{\raggedright\arraybackslash}p{0.22\textwidth}
    >{\raggedright\arraybackslash}X
    >{\raggedright\arraybackslash}X
}
\toprule
\textbf{Aspecto} & \textbf{Progenitor (padre)} & \textbf{Abuela} \\
\midrule
\textbf{Orden} &
Se empieza por él. &
Recién después; la obligación es subsidiaria. \\
\addlinespace
\textbf{Paso previo} &
Mediación contra el progenitor. &
No se avanza hasta que el progenitor no cumpla o se niegue. \\
\addlinespace
\textbf{Si no se presenta} &
Se cierra el acta y queda una constancia. &
Con esa constancia se avanza según lo explicado. \\
\addlinespace
\textbf{Vía} &
Demanda y juicio. &
Después del juicio, se pasa a la abuela. \\
\bottomrule
\end{tabularx}
\end{table}

\subsection*{Cuadro Cornell: caso práctico de alimentos}

\begin{table}[H]
\centering
\begin{tabularx}{\textwidth}{
    >{\raggedright\arraybackslash}p{0.28\textwidth}
    >{\raggedright\arraybackslash}X
}
\toprule
\textbf{Preguntas / palabras clave} &
\textbf{Notas / respuestas} \\
\midrule

\textbf{Caso de alimentos: ¿qué se hace y en qué orden?} &
Se sube primero el escrito de apelación (alimentos provisorios, pedido del 13 de octubre). Luego se busca que se fije la cuota, se demanda la atribución del hogar para que él pueda volver y se plantea el divorcio. \\

\textbf{¿Se puede demandar primero a la abuela?} &
No: la obligación de la abuela es subsidiaria. Se empieza por el progenitor (mediación; si no se presenta se cierra el acta y queda una constancia; demanda y juicio) y después se pasa a la abuela. \\

\midrule

\multicolumn{2}{p{\dimexpr\textwidth-2\tabcolsep\relax}}{
\textbf{Resumen:}
El caso ilustró el orden procesal: primero el escrito de apelación sobre los alimentos provisorios; luego la fijación de la cuota, la demanda por la atribución del hogar y el planteo de divorcio. Ante el reclamo a la abuela, se remarcó su carácter subsidiario: se empieza por el progenitor y solo después se pasa a la abuela.
} \\

\bottomrule
\end{tabularx}
\end{table}

% ============================================================
% SÍNTESIS GENERAL
% ============================================================
\chapter[Síntesis general]{Síntesis general de la clase}

\begin{table}[H]
\centering
\begin{tabularx}{\textwidth}{
    >{\raggedright\arraybackslash}p{0.28\textwidth}
    >{\raggedright\arraybackslash}X
}
\toprule
\textbf{Preguntas / palabras clave} &
\textbf{Notas / respuestas} \\
\midrule

\textbf{¿Qué ciclo completo mostró la clase?} &
El ciclo de un caso en un consultorio jurídico gratuito: derivarlo bien, cerrarlo bien y, mientras está abierto, llevarlo con estrategia. \\

\textbf{¿Cómo se deriva y cómo se cierra?} &
La derivación procede por una causa real, la autoriza la coordinadora, se acredita con certificado o nota firmada y se concreta con un mail <<listo para reenviar>>. Después se renuncia, se pide la desvinculación, se registra la derivación en el drive y se carga la baja administrativa. \\

\textbf{¿Qué ilustró el caso de alimentos?} &
El orden procesal (apelación, fijación de la cuota, atribución del hogar, divorcio) y el carácter subsidiario del reclamo a la abuela. \\

\midrule

\multicolumn{2}{p{\dimexpr\textwidth-2\tabcolsep\relax}}{
\textbf{Resumen:}
La clase mostró el ciclo completo de un caso en un consultorio jurídico gratuito: derivarlo bien, cerrarlo bien y, mientras está abierto, llevarlo con estrategia. Una derivación solo procede por una causa real (típicamente un cambio laboral), la autoriza la coordinadora, se acredita con certificado o nota firmada y se concreta con un mail <<listo para reenviar>> que lleva la declaración jurada, la constancia y la comisión de destino. Después se renuncia, se pide la desvinculación para no quedar notificados, se registra la derivación en el drive y se carga la baja administrativa, lo que ordena el traspaso de fin de cuatrimestre. El caso de alimentos ilustró el orden procesal: primero el escrito de apelación sobre los alimentos provisorios; luego la fijación de la cuota, la demanda por la atribución del hogar y el planteo de divorcio. Finalmente, ante el reclamo a la abuela, el docente remarcó su carácter subsidiario: se empieza por el progenitor (mediación y, si no se presenta, acta con constancia; luego demanda y juicio) y solo después se pasa a la abuela.
} \\

\bottomrule
\end{tabularx}
\end{table}

% ============================================================
% APÉNDICE
% ============================================================
\appendix
\chapter{Material complementario: referencia normativa}

\begin{note}
Este complemento figura en el documento original como referencia editorial, no como parte de lo dicho por el docente. En el minuto 13:12 de la grabación el docente alude a <<el artículo>>, pero el número no se escucha. Como orientación sobre el tema (alimentos entre parientes, obligación subsidiaria de los ascendientes), el documento original transcribe el texto del Código Civil y Comercial de la Nación. \textbf{Verificar contra el texto oficial vigente.}
\end{note}

\textbf{Art. 537.} \textit{<<Los parientes se deben alimentos en el siguiente orden: a) los ascendientes y descendientes. Entre ellos, están obligados preferentemente los más próximos en grado; b) los hermanos bilaterales y unilaterales. En cualquiera de los supuestos, los alimentos son debidos por los que están en mejores condiciones para proporcionarlos. Si dos o más de ellos están en condiciones de hacerlo, están obligados por partes iguales, pero el juez puede fijar cuotas diferentes, según la cuantía de los bienes y cargas familiares de cada obligado.>>}

\textbf{Art. 668 (cita parcial).} \textit{Los alimentos a los ascendientes pueden ser reclamados en el mismo proceso en que se demanda a los progenitores o en proceso diverso; además de lo previsto en el título del parentesco, debe acreditarse verosímilmente las dificultades del actor para percibir los alimentos del progenitor obligado.}

\end{document}
